% !TeX root = ../protokoll.tex
\section{Mechanismus}
\label{sec:mechanismus}

Die Reaktion verläuft nach dem \textbf{Markovnikov-Prinzip}: Das
elektrophile Wasserstoff-Atom lagert sich an das wasserstoffreichere
Kohlenstoffatom an.

\subsection{Schritt 1: Protonierung}
Angriff der Doppelbindung auf das Proton der Säure, Bildung des
tertiären Carbenium-Ions (stabilisiert durch +I-Effekt der
Methylgruppen):
\begin{reaction}
  \ce{(CH3)2C=CHCH3} + \ce{H+} -> \ce{(CH3)2C+ -CH2CH3}
\end{reaction}

\subsection{Schritt 2: Nucleophiler Angriff}
Wasser greift nucleophil am Carbenium-Ion an:
\begin{reaction}
  \ce{(CH3)2C+ -CH2CH3} + \ce{H2O} -> \ce{(CH3)2C(OH2+)CH2CH3}
\end{reaction}

\subsection{Schritt 3: Deprotonierung}
\begin{reaction}
  \ce{(CH3)2C(OH2+)CH2CH3} -> \ce{(CH3)2C(OH)CH2CH3} + \ce{H+}
\end{reaction}

Strukturformel der Zielverbindung (Platzhalter für chemfig):
\begin{figure}[H]
  \centering
  \chemfig{H_3C-C(-[2]OH)(-[3]CH_3)-CH_2-CH_3}
  \caption{2-Methyl-2-butanol.}
  \label{fig:chemfig}
\end{figure}
